\documentclass[11pt,a4paper]{article}


\usepackage[utf8]{inputenc}
\usepackage[T1]{fontenc}
\usepackage[portuguese]{babel}

\usepackage[a4paper,margin=2cm]{geometry}
\usepackage{graphicx}
\usepackage{amsmath}
\usepackage{booktabs}
\usepackage{hyperref}
\usepackage{float}


\hypersetup{
    colorlinks=true,
    linkcolor=black,
    urlcolor=black
}

\setlength{\parskip}{0.5em}
\setlength{\parindent}{0pt}


\begin{document}


\begin{center}

    \Large \textbf{Relatório de Programação Concorrente 25/26} \\[1cm]

    Miguel Gonçalves A108478 \\
    João Taveira A108559 \\
    José Novais A108395 \\
    Luís Lopes A108396 \\[1cm]

    Professor: Paulo Sérgio Soares Almeida \\[1cm]

    \today

\end{center}

\vspace{1cm}
\newpage


\renewcommand{\contentsname}{Índice}
\tableofcontents
\newpage


\section{Introdução}
No âmbito da unidade curricular de Programação Concorrente foi proposto o desenvolvimento de um jogo multijogador em tempo real. O projeto tem como principal objetivo aplicar conceitos relacionados com concorrência, comunicação em rede e sincronização entre processos, recorrendo às linguagens Erlang e Java.
O jogo desenvolvido permite a interação simultânea de vários jogadores num ambiente bidimensional, onde cada utilizador controla um avatar capaz de se movimentar e interagir com outros elementos do jogo.
Ao longo do desenvolvimento foram abordados diferentes desafios associados à gestão de múltiplos clientes, troca contínua de informação e atualização consistente do estado do jogo.
Este relatório apresenta a análise do problema proposto, a solução desenvolvida, a estrutura do projeto e as principais conclusões obtidas durante a implementação.
\newpage



\section{Desenvolvimento}

\subsection{Análise do Problema}
O principal desafio do projeto consiste em garantir o funcionamento correto de um jogo multijogador em tempo real, assegurando que vários jogadores conseguem interagir simultaneamente sem comprometer a consistência do estado do jogo.
O sistema deve gerir autenticação de utilizadores, criação de partidas, movimentação dos jogadores, colisões, capturas e atualização das pontuações. Além disso, é necessário garantir que a informação recebida por cada cliente se mantém sincronizada ao longo da partida.
Relativamente aos requisitos funcionais, o sistema deve permitir:
\begin{itemize}
    \item Registo, autenticação e remoção de utilizadores;
    \item Criação de partidas com 3 a 4 jogadores;
    \item Movimentação dos jogadores num espaço 2D;
    \item Interação entre jogadores e objetos do cenário;
    \item Sistema de capturas e pontuação;
    \item Atualização contínua do estado do jogo.
\end{itemize}
Quanto aos requisitos não funcionais, destacam-se:
\begin{itemize}
    \item Suporte para múltiplas partidas em simultâneo;
    \item Comunicação eficiente entre cliente e servidor;
    \item Capacidade de processamento concorrente;
    \item Sincronização do estado do jogo entre jogadores.
\end{itemize}



\subsection{Solução Implementada}
Para responder aos requisitos do projeto foi adotada uma arquitetura cliente-servidor. O servidor ficou responsável pela execução da lógica do jogo e gestão das partidas, enquanto o cliente trata da interface gráfica e interação do utilizador.
A implementação do servidor em Erlang permitiu utilizar concorrência baseada em processos leves, facilitando a execução simultânea de várias tarefas independentes, como gestão de clientes, matchmaking e partidas ativas.
A comunicação entre cliente e servidor é realizada através de sockets TCP, possibilitando a troca contínua de mensagens relacionadas com ações dos jogadores e atualização do estado do jogo.
O cliente foi desenvolvido em Java com recurso à biblioteca Processing, utilizada para representar graficamente o ambiente de jogo e os diferentes elementos presentes na partida.
Esta abordagem permitiu separar claramente a lógica da aplicação da componente gráfica, tornando o sistema mais organizado e modular.

\newpage 


\section{Implementação}

\subsection{Estrutura do Projeto}
O projeto encontra-se dividido em duas componentes principais: servidor e cliente.
A componente do servidor está organizada em vários módulos independentes, cada um responsável por uma funcionalidade específica do sistema:
\begin{itemize}
    \item \textbf{tcp\_server}: aceita novas ligações TCP;
    \item \textbf{client\_handler}: gere a comunicação individual com cada cliente;
    \item \textbf{matchmaker}: organiza os jogadores em filas de espera e cria partidas;
    \item \textbf{game\_session}: executa a lógica de cada partida;
    \item \textbf{top\_manager}: mantém o registo das pontuações.
\end{itemize}
O cliente encontra-se dividido entre componentes responsáveis pela interface gráfica, processamento do input do utilizador e comunicação com o servidor.
A interface gráfica possui diferentes estados da aplicação, incluindo autenticação, fila de espera e jogo em execução.

\subsection{Comunicação Cliente-Servidor}
A troca de informação entre cliente e servidor é realizada de forma contínua ao longo da partida. Sempre que o jogador executa uma ação de movimento, o cliente envia um comando correspondente ao servidor, indicando as teclas atualmente pressionadas.
Após processar todas as ações dos jogadores, o servidor atualiza o estado da partida e envia para cada cliente informação relativa às posições dos jogadores, objetos presentes no mapa, massas e pontuações atuais.
Para reduzir complexidade na comunicação, foi utilizado um protocolo textual simples, permitindo representar facilmente os diferentes tipos de mensagens trocadas entre as duas componentes do sistema.


\subsection{Gestão Concorrente das Partidas}
Cada partida possui um ciclo de atualização independente, responsável por executar periodicamente a simulação do jogo. Durante cada atualização são calculados movimentos, colisões, capturas e alterações de pontuação.
A separação das partidas em processos distintos permitiu executar vários jogos em simultâneo sem interferência entre sessões. Desta forma, operações realizadas numa partida não afetam o estado das restantes.
A utilização de concorrência em Erlang simplificou também a gestão dos jogadores ligados ao servidor, permitindo tratar múltiplas conexões e eventos em paralelo.


\newpage
\section{Conclusão}

Como conclusão deste trabalho, consideramos que foi uma forma eficaz de aplicar os conteúdos lecionados nas aulas práticas e teóricas, nomeadamente a utilização de barreiras no processo de matchmaking. O resultado final correspondeu às nossas expectativas, refletindo uma implementação sólida dos conceitos abordados.

Adicionalmente, conseguimos implementar concorrência em Java com recurso a threads e mecanismos de sincronização, como locks, o que contribuiu para uma melhor organização e controlo das diferentes operações do sistema, incluindo no Processing, através da utilização de uma thread dedicada a cada tipo de operação. Por fim stamos satisfeitos com o trabalho desenvolvido e com os resultados alcançados.





\end{document}